\chapter{Cardiac Resynchronization Therapy}

\section*{Overview}
Cardiac resynchronization therapy (CRT) uses a specialized pacemaker, with or without a defibrillator, to coordinate contraction of the heart ventricles.

\section*{Why This Test or Procedure Is Ordered}
It may benefit selected patients with symptomatic heart failure, reduced ejection fraction, and electrical conduction delay such as a wide QRS complex.

\section*{How This Test or Procedure Is Performed}
A pulse generator is placed beneath the skin, usually below the collarbone. Leads are positioned in the right atrium or ventricle and through a cardiac vein on the left ventricle, then programmed to coordinate pacing.

\section*{Risks and Follow-up}
Infection, bleeding, pneumothorax, lead displacement, vein injury, arrhythmia, and device malfunction are possible. Regular device checks and heart-failure treatment remain necessary.

\section*{References}
\begin{enumerate}
  \item MedlinePlus. Biventricular Pacemaker. U.S. National Library of Medicine; 2025.
  \item Current Medical Diagnosis and Treatment. McGraw-Hill; 2025.
\end{enumerate}
\clearpage
